\documentclass[a4paper,12pt]{article}
\usepackage[T2A]{fontenc}
\usepackage[utf8]{inputenc}
\usepackage[russian]{babel}
\usepackage{amsmath,amssymb,mathtools}
\usepackage{geometry}
\geometry{left=22mm,right=22mm,top=23mm,bottom=25mm}
\usepackage{microtype}
\renewcommand{\familydefault}{\sfdefault}
\usepackage{booktabs,tabularx,array,enumitem}
\usepackage{tikz}
\usepackage{xcolor}
\usepackage{fancyhdr}
\usepackage{hyperref}
\usepackage{parskip}
\definecolor{accent}{HTML}{006B70}
\definecolor{textgray}{HTML}{40474B}
\definecolor{errorred}{HTML}{A42431}
\hypersetup{colorlinks=true,linkcolor=accent,urlcolor=accent}
\setlength{\parindent}{0pt}
\setlength{\parskip}{6pt}
\setlist[enumerate]{leftmargin=*,itemsep=3pt,topsep=3pt}
\setlength{\emergencystretch}{2em}
\pagestyle{fancy}
\fancyhf{}
\renewcommand{\headrulewidth}{0pt}
\fancyfoot[C]{\footnotesize\color{textgray}Лёвин Артём Александрович эксклюзивно для Марка Гукина}
\fancyfoot[R]{\footnotesize\thepage}
\newcommand{\rubric}[1]{\par\medskip\textbf{#1}\par\nobreak}
\newcommand{\err}[1]{\textcolor{errorred}{\textbf{Ошибка:}} #1}
\newcommand{\keyidea}[1]{\textcolor{accent}{\textbf{Ключевая идея:}} #1}
\newcommand{\outcome}[1]{\textcolor{accent}{\textbf{Итог:}} #1}
\begin{document}
\begin{titlepage}
\thispagestyle{empty}
\centering
\vspace*{30mm}
{\large\color{accent} МАТЕМАТИКА\par}
\vspace{21mm}
{\Huge\bfseries Ревью домашней работы\par}
\vspace{13mm}
{\Large Тригонометрия. Единичная окружность,\\
табличные значения и уравнения\par}
\vfill
{\large 26 октября 2026 г.\par}
\vspace{9mm}
{\large Лёвин Артём Александрович\par}
\vspace{4mm}
{\normalsize эксклюзивно для Марка Гукина\par}
\vspace{17mm}
\begin{tikzpicture}
\draw[accent!55,line width=0.8pt] (-5,0) .. controls (-2,0.75) and (0,-0.65) .. (1.7,-0.1) .. controls (3,0.25) and (4,-0.15) .. (5,0.12);
\end{tikzpicture}
\end{titlepage}

\section*{Сводка замечаний}
В таблицу включены только задания, где требуется исправление. Номер задания ведёт к подробному разбору.
\par\vspace{4mm}
\noindent{\small
\renewcommand{\arraystretch}{1.4}
\begin{tabularx}{\linewidth}{@{}p{0.8cm}>{\raggedright\arraybackslash}p{2.0cm}>{\raggedright\arraybackslash}X p{3.5cm}p{3.9cm}@{}}
\toprule
\textbf{№} & \textbf{Тема} & \textbf{Суть ошибки} & \textbf{Ваш ответ} & \textbf{Правильный ответ}\\
\midrule
\hyperref[task:1]{\textbf{1}} & Единичная окружность & Приближённые координаты вместо точных; неверный знак абсциссы для $-45^\circ$. & $45^\circ:(0{,}6;\,0{,}6)$; $-45^\circ:(-0{,}6;\,-0{,}6)$ & $45^\circ:\left(\frac{\sqrt2}{2};\frac{\sqrt2}{2}\right)$; $-45^\circ:\left(\frac{\sqrt2}{2};-\frac{\sqrt2}{2}\right)$ \\
\addlinespace
\hyperref[task:4]{\textbf{4}} & Табличные значения & Косинус $60^\circ$ перепутан со значением для $45^\circ$. & $\cos\frac{\pi}{3}=\frac{\sqrt2}{2}$ & $\cos\frac{\pi}{3}=\frac12$ \\
\addlinespace
\hyperref[task:6]{\textbf{6а}} & Уравнение & Смешаны градусы и радианы в общей формуле корней. & $\alpha=270+2\pi n$ & $\alpha=\frac{3\pi}{2}+2\pi n$ \\
\bottomrule
\end{tabularx}
}\par
\vspace{6mm}
\noindent\textcolor{accent}{\textbf{Порядок работы:}} прочитайте объяснение первой ошибки, восстановите ход решения самостоятельно и только затем переходите к задаче на закрепление.

\clearpage
\thispagestyle{fancy}
\vspace*{\fill}
\begin{center}
{\LARGE\bfseries Разбор ошибок}\par
\vspace{7mm}
{\large Тригонометрия}\par
\end{center}
\vspace*{\fill}
\clearpage

\phantomsection\label{task:1}
\section*{Задание 1. Точки на единичной окружности}
\rubric{Условие}
На одной единичной окружности отметьте точки для углов $45^\circ$, $-45^\circ$, $180^\circ$, $210^\circ$, $270^\circ$. Запишите \textbf{точные} координаты каждой точки.

\rubric{Ваше решение}
В записи координат указаны, в частности,
\[
45^\circ\mapsto(0{,}6;\,0{,}6),\quad
-45^\circ\mapsto(-0{,}6;\,-0{,}6),\quad
210^\circ\mapsto(-0{,}9;\,-0{,}3).
\]
Точки $180^\circ$ и $270^\circ$ записаны как $(-1;0)$ и $(0;-1)$.

\rubric{Ваш ответ}
Для трёх углов приведены десятичные координаты; для $-45^\circ$ обе координаты указаны отрицательными.

\rubric{Ошибка}
\err{Координаты точек на единичной окружности необходимо записывать точно. У точки, соответствующей $-45^\circ$, абсцисса положительна.}

\rubric{Где именно возникает ошибка}
Верно определено, что угол $45^\circ$ лежит в первой четверти и обе его координаты положительны. Первый неверный шаг~--- запись $(0{,}6;\,0{,}6)$ вместо точных значений. При переходе к углу $-45^\circ$ допущена дополнительная ошибка: точка расположена в четвёртой четверти, но первой координате приписан знак минус.

\rubric{Почему это неверно}
На единичной окружности любая точка $(x;y)$ удовлетворяет равенству $x^2+y^2=1$. Для пары $(0{,}6;0{,}6)$ получаем $0{,}36+0{,}36=0{,}72$, поэтому такая точка даже не лежит на единичной окружности. Аналогично для записи $(-0{,}9;-0{,}3)$ сумма квадратов равна $0{,}81+0{,}09=0{,}9$.

Кроме того, отрицательный угол отсчитывается от положительного направления оси $x$ \emph{по часовой стрелке}. Поэтому $-45^\circ$ приводит в четвёртую четверть: $x>0$, $y<0$.

\rubric{Ключевая идея}
\keyidea{На единичной окружности координаты точки равны $x=\cos\alpha$, $y=\sin\alpha$. Сначала определите четверть и знаки, затем найдите точные табличные значения.}

\rubric{Исправление}
Для угла $45^\circ$ обе координаты равны и положительны: если $x=y$, то $2x^2=1$ и $x=y=\frac{\sqrt2}{2}$. Угол $-45^\circ$ имеет ту же положительную абсциссу и противоположную ординату. Для угла $210^\circ$ используйте $210^\circ=180^\circ+30^\circ$.

\rubric{Полное правильное решение}
\begin{enumerate}
\item $45^\circ$~--- первая четверть. $\cos45^\circ=\sin45^\circ=\frac{\sqrt2}{2}$, значит
\[P_{45}=\left(\frac{\sqrt2}{2};\,\frac{\sqrt2}{2}\right).\]
\item $-45^\circ$~--- четвёртая четверть. Косинус положительный, синус отрицательный:
\[P_{-45}=\left(\frac{\sqrt2}{2};\,-\frac{\sqrt2}{2}\right).\]
\item $180^\circ$~--- отрицательная полуось $x$:
\[P_{180}=(-1;\,0).\]
\item $210^\circ=180^\circ+30^\circ$~--- третья четверть. Опорный угол $30^\circ$, обе координаты отрицательные:
\[P_{210}=\left(-\frac{\sqrt3}{2};\,-\frac12\right).\]
\item $270^\circ$~--- отрицательная полуось $y$:
\[P_{270}=(0;\,-1).\]
\end{enumerate}

\rubric{Проверка}
Например, для $45^\circ$ сумма квадратов координат равна $\frac24+\frac24=1$; для $210^\circ$~--- $\frac34+\frac14=1$. Знаки также соответствуют четвертям. Для точек на осях сумма квадратов равна единице непосредственно.

\rubric{Итог}
\outcome{Точные координаты задаются через корни и дроби. Для $-45^\circ$ первая координата обязательно положительна; для $210^\circ$ обе отрицательны.}

\rubric{Задача на закрепление}
Отметьте на единичной окружности точки для углов $30^\circ$, $-60^\circ$, $150^\circ$, $225^\circ$ и $300^\circ$. Запишите точные координаты каждой точки.

\clearpage
\phantomsection\label{task:4}
\section*{Задание 4. Табличные значения}
\rubric{Условие}
Восстановите значения $\sin\frac{\pi}{6}$, $\cos\frac{\pi}{3}$, $\sin\frac{\pi}{4}$ и $\cos\frac{\pi}{6}$.

\rubric{Ваше решение}
В записи приведён переход
\[\cos\frac{\pi}{3}=\cos60^\circ=\frac{\sqrt2}{2}.\]
\rubric{Ваш ответ}
$\displaystyle\cos\frac{\pi}{3}=\frac{\sqrt2}{2}$.

\rubric{Ошибка}
\err{Число $\frac{\sqrt2}{2}$ соответствует косинусу $45^\circ$, а косинус $60^\circ$ равен $\frac12$.}

\rubric{Где именно возникает ошибка}
Переход от радиан к градусам верен:
\[\frac{\pi}{3}=60^\circ.\]
Первый неверный шаг~--- подстановка табличного значения $\frac{\sqrt2}{2}$ после правильной записи $\cos60^\circ$.

\rubric{Почему это неверно}
Углы $30^\circ$ и $60^\circ$ дополняют друг друга до $90^\circ$. Поэтому косинус $60^\circ$ равен синусу $30^\circ$. В прямоугольном треугольнике с углами $30^\circ$, $60^\circ$, $90^\circ$ катет против угла $30^\circ$ составляет половину гипотенузы. Значит,
\[\cos60^\circ=\sin30^\circ=\frac12.\]
Значение $\frac{\sqrt2}{2}$ возникает для угла $45^\circ$ в равнобедренном прямоугольном треугольнике.

\rubric{Ключевая идея}
\keyidea{Сначала переведите радианный угол в градусы, затем обратитесь к нужному углу в таблице $30^\circ$, $45^\circ$, $60^\circ$. Не смешивайте соседние столбцы.}

\rubric{Исправление}
В последнем шаге заменяем значение функции:
\[\cos\frac{\pi}{3}=\cos60^\circ=\boxed{\frac12}.\]

\rubric{Полное правильное решение}
\begin{align*}
\sin\frac{\pi}{6}&=\sin30^\circ=\frac12,\\
\cos\frac{\pi}{3}&=\cos60^\circ=\frac12,\\
\sin\frac{\pi}{4}&=\sin45^\circ=\frac{\sqrt2}{2},\\
\cos\frac{\pi}{6}&=\cos30^\circ=\frac{\sqrt3}{2}.
\end{align*}

\rubric{Проверка}
Основное тригонометрическое тождество при $60^\circ$ даёт
\[\sin^260^\circ+\cos^260^\circ
=\frac34+\frac14=1.\]
Если вместо $\frac12$ ошибочно поставить $\frac{\sqrt2}{2}$, сумма станет $\frac34+\frac24=\frac54$, что противоречит тождеству.

\rubric{Итог}
\outcome{$\cos60^\circ=\frac12$; $\cos45^\circ=\frac{\sqrt2}{2}$; $\cos30^\circ=\frac{\sqrt3}{2}$.}

\rubric{Задача на закрепление}
Вычислите точно $\displaystyle\cos\frac{2\pi}{3}$.

\clearpage
\phantomsection\label{task:6}
\section*{Задание 6а. Уравнение $\sin\alpha=-1$}
\rubric{Условие}
Решите уравнение $\sin\alpha=-1$ и запишите все его решения.

\rubric{Ваше решение}
В тетради указана запись вида
\[\alpha=270+2\pi n,\qquad n\in\mathbb Z.\]
\rubric{Ваш ответ}
$\alpha=270+2\pi n$.

\rubric{Ошибка}
\err{В одной формуле объединены число $270$, соответствующее градусному углу, и период $2\pi$, записанный в радианах.}

\rubric{Где именно возникает ошибка}
Верно найдено начальное направление на единичной окружности: синус равен $-1$ в нижней точке, соответствующей углу $270^\circ$. Ошибка возникает при записи серии решений: угол дан числом $270$, а прибавляемый период взят как $2\pi$.

\rubric{Почему это неверно}
Период полного оборота выражается либо в градусах ($360^\circ$), либо в радианах ($2\pi$). Запись одной формулы должна использовать \emph{одну единицу измерения}. Число $270$ без значка градуса при радианной записи означает $270$ радиан; это совершенно другой угол. Если иметь в виду $270^\circ$, к нему следует прибавлять $360^\circ n$. При выборе радиан $270^\circ$ нужно заменить на $\frac{3\pi}{2}$.

\rubric{Ключевая идея}
\keyidea{Сначала найдите точку на окружности, затем запишите один подходящий угол и прибавьте целое число полных оборотов в тех же единицах.}

\rubric{Исправление}
Правильная градусная запись:
\[\alpha=270^\circ+360^\circ n,\qquad n\in\mathbb Z.\]
Поскольку $270^\circ=\frac{3\pi}{2}$, в радианах получаем
\[\alpha=\frac{3\pi}{2}+2\pi n,\qquad n\in\mathbb Z.\]

\rubric{Полное правильное решение}
\begin{enumerate}
\item На единичной окружности синус~--- это ордината точки.
\item Наименьшая ордината равна $-1$ и достигается только в точке $(0;-1)$.
\item Один угол, соответствующий этой точке, равен $\frac{3\pi}{2}$; другой удобный представитель~--- $-\frac{\pi}{2}$.
\item Каждое прибавление или вычитание полного оборота $2\pi$ возвращает в ту же точку.
\end{enumerate}
Следовательно,
\[\boxed{\alpha=\frac{3\pi}{2}+2\pi n,\quad n\in\mathbb Z.}\]
Равносильная запись: $\alpha=-\frac{\pi}{2}+2\pi n$, где $n\in\mathbb Z$.

\rubric{Проверка}
\[\sin\left(\frac{3\pi}{2}+2\pi n\right)
=\sin\frac{3\pi}{2}=-1\]
для любого целого $n$. Иных точек с ординатой $-1$ на единичной окружности нет.

\rubric{Итог}
\outcome{Общая формула решения в радианах: $\alpha=\frac{3\pi}{2}+2\pi n$, $n\in\mathbb Z$.}

\rubric{Задача на закрепление}
Решите уравнение $\sin\alpha=1$. Запишите все решения в радианах.
\end{document}